\documentclass[10pt,a4paper]{amsart}
\usepackage[margin=19mm]{geometry}
\usepackage{amsmath,amssymb,mathtools}
\usepackage[scr=rsfs,cal=euler]{mathalfa}
\usepackage{xfrac}
\usepackage[unicode,hidelinks,bookmarks=false]{hyperref}
\newcommand{\ie}{i.e.\ }
\newcommand{\cf}{cf.\ }
\newcommand{\resp}{respectively\ }
\newcommand{\Subsection}{\subsection}
\providecommand{\nobreakdash}{\nobreak\hbox{-}}
\setlength{\emergencystretch}{2em}
\multlinegap0pt
\usepackage{amssymb}
\usepackage{xcolor}
\newtheorem{theorem}{Theorem}
\newtheorem{corollary}{Corollary}
\newtheorem{lemma}{Lemma}
\def\Prob{{\mathbb P}}
\def\Zmax{Z_{(n)}}
\title{Non-standard boundary behaviour in two-component mixture models}
\author{Heather Battey, \quad Peter McCullagh \quad and \quad Daniel Xiang$^{\dagger}$}
\date{Source: arXiv:2407.20162v3 (CC BY 4.0)}
\begin{document}
\maketitle

\noindent\emph{Source and licence.} Non-standard boundary behaviour in two-component mixture models. Version arXiv:2407.20162v3 (CC BY 4.0), distributed by arXiv under the Creative Commons Attribution 4.0 International licence, http://creativecommons.org/licenses/by/4.0/, as recorded in the arXiv licence metadata for this exact version and shown on the abstract page https://arxiv.org/abs/2407.20162v3. Full licence notice: \texttt{licenses/arxiv-2407.20162-CC-BY-4.0.txt}.\par\smallskip

\noindent\textit{Editorial scope.} Counted unit: the article's lemma lem:sum-max-dependence (source0.tex lines 718--730). The article's complete original proof of it is delivered with it (source0.tex lines 732--772, one \texttt{proof} environment of 41 lines). No mathematics is written, repaired or re-derived: the statement and the proof are the article's own lines, in the article's own notation, and no retained line is substituted or re-flowed.  Retained uncounted as the source-local context the proof needs: lines 349--351; lines 397--399; lines 402--417; lines 452--461; lines 465--476; lines 692--702.  Nothing was omitted from any retained span, so the package carries no omission record.  No retained line contains a citation, so the wrapper carries no bibliography and no bibliography entry.  No retained line contains a figure environment, an image-inclusion command or drawing code, so no image is delivered and the package declares no asset dependency.  Editorial formatting only, in addition: an AMS \texttt{amsart} preamble that restates the article's own declarations and macros used by the retained lines, namely the environments \texttt{theorem}, \texttt{corollary}, \texttt{lemma} on separate counters, together with the standard packages those lines need.  The retained environments print in this document as Theorem 1, Corollary 1, Lemma 1 in order of appearance, whereas the article prints the same items with the numbering of its own full document.  The complete native source of the article as distributed in the arXiv e-print for this exact version is bundled unchanged as \texttt{sources/163054/source0.tex}.
\begin{equation}\label{eqDB}
L(x L^\dagger(x)) L^\dagger(x) \sim L^\dagger(x L(x)) L(x) \sim 1,
\end{equation}

\begin{equation}\label{eqAn}
A_n \sim n \Im \log \psi(1/B_n) \sim n \int \sin(x/B_n) \, dF(x).
\end{equation}

\begin{theorem}\label{thmSinIntegral}
	Let $F$ be a finite-mean distribution on the positive real line whose tail is $\bar F(x) \sim 2C_1/(x L(x))$, where $C_1=1/\pi$ and the slowly varying function $L(x)$ admits the parametrization
	%
	\begin{equation}\label{eqLParam}
	L(x) = \left\{ \begin{array}{ll}(\beta_0\log x)^{\delta + 1} e^{(\beta_1\log x)^\gamma} & \beta_1> 0 \\
	(\beta_0\log x)^{\delta + 1} & \beta_1 = 0,
	\end{array}\right.
	\end{equation}
	%
	for $\beta_0>0$. Finiteness of the mean implies either $\beta_1>0$ and $0 < \gamma < 1$,	or $\beta_1 = 0$ and $\delta > 0$. The sine-integral for large $T = 1/|t|$ is,
	\begin{eqnarray*}
		\int_0^\infty \sin (t x)  \, dF(x) 
		&=& \mu t - K_{\delta,\gamma,\beta_1}  \frac{t\,(\log T)^{1-\gamma}} {L(T)} +o\biggl(\frac {t (\log T)^{1-\gamma}} {L(T)} \biggr), \quad |t|\rightarrow 0,
	\end{eqnarray*}
	where $\mu$ is the mean, $K_{\delta,\gamma, \beta_1} = 2C_1/ (\beta_1^{\gamma}\gamma)$ for $\gamma, \beta_1 > 0$, and $K_{\delta,0,0} = 2C_1/\delta$.
\end{theorem}

\begin{corollary}\label{corollAn}
	For a distribution satisfying the conditions of Theorem \ref{thmSinIntegral}, the centering sequence \eqref{eqAn} is 
	%
	\begin{eqnarray*}
		A_n &=& \frac {n \mu} {B_n} - K_{\delta,\gamma,\beta_1}  \frac{n (\log B_n)^{1-\gamma}}  {B_n L(B_n)} (1 + o(1)),\\
		&=& \frac {\mu} {L^\dagger(n)} - K_{\delta,\gamma,\beta_1}  \frac{(\log B_n)^{1-\gamma}} {L^\dagger(n) L(n L^\dagger(n))}(1 + o(1)), \\
		&=& \frac {\mu} {L^\dagger(n)} - K_{\delta,\gamma,\beta_1}(\log n)^{1-\gamma} (1 + o(1)).
	\end{eqnarray*}
	The third line follows from the definition of the de~Bruijn conjugate in \eqref{eqDB}.
\end{corollary}

\begin{corollary}\label{corollXbar}
	For a distribution satisfying the conditions of Theorem \ref{thmSinIntegral}, the distribution of the sample average for large $n$ is
	\[
	\frac{\bar X_n - \mu} {L^\dagger(n)} = - K_{\delta,\gamma,\beta_1}  (\log n)^{1-\gamma} + \varepsilon + o_p(1),
	\]
	where $\varepsilon$ has the Cauchy limit with skewness $\beta=1$.
	Given the right-tail behaviour of the limit distribution,
	\[
	\Prob(\bar X_n > \mu) \sim 2 K_{\delta,\gamma,\beta_1}^{-1} (\log n)^{\gamma-1} / \pi, 
	\]
	which tends to zero at rate $\beta_1^\gamma \gamma (\log n)^{\gamma-1}$ if $\beta_1,\gamma > 0$, or $\delta/\log n$ if $\beta_1=0$.
\end{corollary}

\begin{theorem}\label{cjl_theorem}
	Let $F$ be a zero-mean distribution with support $(-1, \infty)$ and tail satisfying the conditions of Theorem~\ref{thmSinIntegral}, and let $Z_1,\ldots, Z_n$ be an iid sample from $F$. Let $T_n = B_n|A_n| = K_{\delta,\gamma,\beta_1} B_n (\log n)^{1-\gamma}$, where $B_n=n/L(n)$ and $A_n$ is given in Corollary~\ref{corollAn}.
	Given $\bar Z_n > 0$, the conditional limit distribution of $(\bar Z_n, \Zmax)$ is such that
	\[
	\biggl(\frac{n\bar Z_n} {T_n},\, \frac{\Zmax} {T_n} \biggr) \sim \biggl(\frac U{1-U},\, \frac 1 {1-U} \biggr),
	\]
	where $U$ is uniform on $(0, 1)$.
	It follows that $n\bar Z_n/T_n \sim T_n / (n\bar Z_n)$ is conditionally self-reciprocal,
	the ratio $n\bar Z_n / \Zmax \sim U$ is conditionally uniform, 
	and the support of the joint distribution degenerates to the line $\Zmax/T_n - n\bar Z_n/T_n -1  \to 0$.
\end{theorem}

\begin{lemma}
	\label{lem:sum-max-dependence}
	Let $F$ be a zero-mean distribution with support $(-1, \infty)$ and tail satisfying the conditions of Theorem~\ref{thmSinIntegral}, and let $Z_1,\ldots, Z_n$ be an iid sample from $F$. Let $B_n,T_n$ be defined as in Theorem \ref{cjl_theorem}. For any fixed $y>0$, conditional on $Z_{(n)}>y T_n$, we have
	\begin{align}
	\label{eq:sum-max-dependence}
	\frac{n\bar{Z}_{n}}{T_n} = \frac{Z_{(n)}+\sum_{k\geq 1}Z_{(n-k)}}{T_n} = \frac{Z_{(n)}}{T_n}-1+o_p(1).
	\end{align}
	Consequently, for any fixed $y>1$
	\begin{align*}
	\lim_{n\to\infty} \; &\Prob\left( \bar{Z}_n>0 \mid Z_{(n)}>yT_n \right) = 1,
	\end{align*}
	so that $\Prob(\bar{Z}_n>0, Z_{(n)}>yT_n) \sim \Prob(Z_{(n)}>yT_n)$ as $n\to\infty$.
\end{lemma}

\begin{proof}
	Let $Z_{(n-k)}$ be the $k$th order statistic with $Z_{(n+1)} = \infty$.
	Since the transformed variables $F(Z_1), \ldots, F(Z_n)$ are independent uniform, the joint distribution of the top order statistics is such that successive differences of the probability-integral transformed values have exactly the same distribution as successive spacings of the top~$k+1$ of $n$ independent uniform order statistics, i.e.
	\[
	n \{\bar F(Z_{(n-r)}) - \bar F(Z_{(n-r+1)}) \}_{0\le r \le k}
	\; \stackrel{(d)}{\to} \, (e_0, \ldots, e_k),
	\]
     where $e_0, \ldots, e_k$ are independent unit exponential variables, whose distribution approximates that of the spacings among the top $k+1$ uniform order statistics, for large $n$ and $k\ll n$. Since $u:=~\bar F(z) \sim  2C_1/(z L(z))$ for large $z$ implies $zL(z)\sim 2C_1/u$, and since the property $L^\dagger(z L (z))L(z) \sim 1$ implies $L(z)\sim 1/L^\dagger(2C_1/u)$, the inverse function inherits the asymptotic behaviour
	\[
	\bar{F}^{-1}(u)\sim (2C_1/u)\; L^\dagger(2C_1/u), \quad  (u \rightarrow 0).
	\] 
	Now since $n\bar{F}(Z_{(n)}) \stackrel{(d)}{\to} e_0$ and 
	\[
	\bar F^{-1}(e_0/n) \;\sim (2C_1n/e_0) \;  L^\dagger(2C_1n/e_0)  \sim \frac{2C_1}{e_0} \; n L^\dagger(n) \sim 2C_1 B_n/e_0,
	\]
	it follows that $Z_{(n)}/B_n \stackrel{(d)}{\to}2C_1/e_0$. 
	Similarly, since $n\bar{F}(Z_{(n-k)}) \stackrel{(d)}{\to} e_0+\dots+e_k$, we obtain
	%
	\begin{equation}\label{e_rep}
	\bar{F}^{-1}((e_0+\dots+e_k)/n) \sim \frac{2C_1 B_n}{e_1+\dots+e_k}
	\end{equation}
	which implies $Z_{(n-k)}/B_n \stackrel{(d)}{\to} 2C_1/(e_0+\dots+e_k)$. 
	
	
	If $\Zmax > yT_n$ for some $y > 0$, then $e_0$~is unusually small compared with $e_1, e_2,\ldots$, so that $e_0 + \cdots + e_k \sim e_1 + \cdots + e_k$. This is because $e_1,\dots,e_k$ are unit exponential variables that are independent of $e_0$, and on the event $\{Z_{(n)}>yT_n\}$, we have
    \begin{align*}
        e_0 \sim n\bar{F}(Z_{(n)}) \leq n\bar{F}(yT_n) \sim \frac{2nC_1}{yT_n L(yT_n)} = \frac{2C_1 L(n)}{K_{\delta,\gamma,\beta_1} yL(yT_n)(\log n)^{1-\gamma}}
    \end{align*}
    by definition of $T_n=K_{\delta,\gamma,\beta_1} n (\log n)^{1-\gamma}/L(n)$. Since $L$ is slowly varying, $\frac{L(yT_n)}{L(n)} \sim 1$ as $n\to\infty$, so the right hand side of the above tends to zero.
    It follows that the conditional distribution of $Z_{(n-1)}, Z_{(n-2)},\ldots$ given $\Zmax > y T_n$ is approximately the same
	as the unconditional distribution of $\Zmax, Z_{(n-1)},\dots$.
	Since the top order statistics $Z_{(n-1)}, Z_{(n-2)},\ldots$ dominate the partial sum $n \bar Z_n - \Zmax$, on the event $\Zmax > y T_n$, Corollary \ref{corollXbar} applied with $\sum_{j \ge1} Z_{(n-j)}$ in place of $n \bar X_n$ implies
	\begin{equation}\label{eq16}
	n \bar Z_n = \Zmax + \sum_{j \ge1} Z_{(n-j)} \sim \Zmax - T_n + O_p(B_n),
	\end{equation}
    %
	which implies \eqref{eq:sum-max-dependence}. For $y>1$, $B_n = o(T_n)$ implies
	\begin{align*}
	\Prob(n\bar{Z}_n > 0 \mid Z_{(n)}>yT_n) \to 1.
	\end{align*}
\end{proof}

\end{document}
